\chaptertitle{Sloth Canon}

                                    Sloth Canon

                    This time, we find Achilles and the Tortoise visiting
                         the dwelling of their new friend, the Sloth.

Achilles: Shall I tell you of my droll footrace with Mr. T?
Sloth: Please do.
Achilles: It has become quite celebrated in these parts. I believe it's even been written up,
   by Zeno.
Sloth: It sounds very exciting.
Achilles: It was. You see, Mr. T began way ahead of me. He had such a huge head start,
   and yet
Sloth: You caught up, didn't you?
Achilles: Yes-being so fleet of foot, I diminished the distance between us at a constant
   rate, and soon overtook him.
Sloth: The gap kept getting shorter and shorter, so you could.
Achilles: Exactly. Oh, look-Mr. T has brought his violin. May I try playing on it, Mr. T?
Tortoise: Please don't. It sounds very flat.
Achilles: Oh, all right. But I'm in a mood for music. I don't know why. Sloth: You can
   play the piano, Achilles.
Achilles: Thank you. I'll try it In a moment. I just wanted to add that I also had another
   kind of "race" with Mr. T at a later date. Unfortunately, in that race
Tortoise: You didn't catch up, did you? The gap kept getting longer and longer, so you
   couldn't.
Achilles: That's true. I believe THAT race has been written up, too, by Lewis Carroll.
   Now, Mr. Sloth, I'll take up your offer of trying out the piano. But I'm so bad at the
   piano. I'm not sure I dare. Sloth: You should try.

   (Achilles sits down and starts playing a simple tune.)

Achilles: Oh-it sounds very strange. That's not how it's supposed to sound at all!
  Something is very wrong.
Tortoise: You can't play the piano, Achilles. You shouldn't try.
Achilles: It's like a piano in a mirror. The high notes are on the left, and the low notes are
  on the right. Every melody comes out inverted, as if upside down. Who would have
  ever thought up something so cockeyed as that?
Tortoise: That's so characteristic of sloths. They hang from
Achilles: Yes, I know-from tree branches-upside down, of course. That sloth-piano would
  he appropriate for playing inverted melodies such




Sloth Canon                                                                               678

FIGURE 133. "Sloth Canon",from the Musical Offering, by J. S. Bach. [Music printed
by Donald Byrd's program "SMUT'.


Sloth Canon                                                                      679

   as occur in some canons and fugues. But to learn to play a piano while hanging from a
   tree must he very difficult. You must have to devote a great deal of energy to it.
Sloth: That's not so characteristic of sloths.
Achilles: No, I gather sloths like to take life very easy. They do everything about half as
   fast as normal. And upside down, to boot. What a peculiar way to go through life!
   Speaking of things that are both upside- and slowed-down, there's a "Canon per
   augmentationem, contrario motu" in the Musical Offering. In my edition, the letters
   `S', `A', `T' are in front of the three staves. I don't know why. Anyway, I think Bach
   carried it off very skillfully. What's your opinion, Mr. T?
Tortoise: He outdid himself. As for those letters "SAT", you could guess what they stand
   for.
Achilles: "Soprano", "Alto", and "Tenor", I suppose. Three-part pieces are often written
   for that combination of voices. Wouldn't you agree, Mr. Sloth?
Sloth: They stand for
Achilles: Oh, just a moment, Mr. Sloth. Mr. Tortoise-why are you putting on your coat?
   You're not leaving, are you? We were just going to fix a snack to eat. You look very
   tired. How do you feel?
Tortoise: Out of gas. So long! (Trudges wearily out the door.)
Achilles: The poor fellow-he certainly looked exhausted. He was jogging all morning.
   He's in training for another race with me. Sloth: He did himself in.
Achilles: Yes, but in vain. Maybe he could beat a Sloth ... but me? Never! Now-weren't
   you about to tell me what those letters "SAT" stand for? Sloth: As for those letters
   "SAT", you could never guess what they stand for.
Achilles: Well, if they don't stand for what I thought, then my curiosity is piqued.
   Perhaps I'll think a little more about it. Say, how do you cook French fries? Sloth: In
   oil.
Achilles: Oh, yes-I remember. I'll cut up this potato into strips an inch or two in length.
Sloth: So short?
Achilles: All right, already, I'll cut four-inch strips. Oh, boy, are these going to be good
   French fries! Too bad Mr. T won't be here to share them.




Sloth Canon                                                                             680

         CHAPTER XX                            modify the subsidiary intentions
                                               but will also modify the rules
                                               which are used in their derivation,
     Strange Loops,                            or in which it will modify the ways
 Or Tangled Hierarchies                        in which it modifies the rules, and
                                               so on, or even in which one
                                               machine will design and construct
      Can Machines Possess
                                               a second machine with enhanced
          Originality?                         capabilities. However, and this is
                                               important, the machine will not
IN THE CHAPTER before last, I                  and cannot [italics are his do any
described       Arthur     Samuel's   very     of these things until it has been
successful checkers program-the one            instructed as to how to proceed.
which can beat its designer. In light of       There is and logically there must
that, it is interesting to hear how Samuel     always remain a complete hiatus
himself feels about the issue of               between (i) any ultimate extension
computers and originality. The following       and elaboration in this process of
extracts are taken from a rebuttal by          carrying out man's wishes and (ii)
Samuel, written in 1960, to an article by      the development within the
Norbert Wiener.                                machine of a will of' its own. To
                                               believe otherwise is either to
   It is my conviction that machines           believe in magic or to believe that
   cannot possess originality in the           the existence of man's will is an
   sense implied by Wiener in his              illusion and that man's actions are
   thesis that "machines can and do            as mechanical as the machine's.
   transcend some of the limitations           Perhaps Wiener's article and my
   of their designers, and that in             rebuttal     have     both     been
   doing so they may be both                   mechanically determined, but this
   effective and dangerous." .. .              I refuse to believe.'
   A machine is not a genie, it does
   not work by magic, it does not            This reminds me of the Lewis Carroll
   possess a will, and, Wiener to the        Dialogue (the Two-Part Invention); I'll
   contrary, nothing comes out which         try to explain why. Samuel bases his
   has not been put in, barring, of          argument against machine consciousness
   course, an infrequent case of             (or will) on the notion that any
   malfunctioning... .                       mechanical instantiation of will would
   The "intentions" which the                require an infinite regress. Similarly,
   machine seems to manifest are the         Carroll's Tortoise argues that no step of
   intentions     of     the     human       reasoning, no matter how simple, can be
   programmer, as specified in               done without invoking some rule on a
   advance, or they are subsidiary           higher level to justify the step in
   intentions derived from these,            question. But that being
   following rules specified by the
   programmer. We can even
   anticipate    higher     levels   of
   abstraction, just as Wiener does, in
   which the program will not only

Strange Loops, Or Tangled Hierarchies                                             681

also a step of reasoning. one must resort    anything without having a rule telling it
to a yet higher-level rule, and so on.       to do so. In fact, machines get around the
Conclusion: Reasoning involves an            Tortoise's silly objections as easily as
infinite regress.                            people do, and moreover for exactly the
         Of course something is wrong        same reason: both machines and people
with the Tortoise's argument, and I          are made of hardware which runs all by
believe something analogous is wrong         itself, according to the laws of physics.
with Samuel's argument. To show how          There is no need to rely on "rules that
the fallacies are analogous, I now shall     permit you to apply the rules", because
"help     the     Devil",   by    arguing    the lowest-level rules-those without any
momentarily as Devil's advocate. (Since,     "meta"'s in front-are embedded in the
as is well known, God helps those who        hardware, and they run without
help themselves, presumably the Devil        permission. Moral: The Carroll Dialogue
helps all those, and only those, who don't   doesn't say anything about the
help themselves. Does the Devil help         differences     between     people    and
himself?) Here are my devilish               machines, after all. (And indeed,
conclusions drawn from the Carroll           reasoning is mechanizable.)
Dialogue:                                             So much for the Carroll
                                             Dialogue. On to Samuel's argument.
   The conclusion "reasoning is              Samuel's point, if I may caricature it, is
   impossible" does not apply to             this:
   people, because as is plain to
   anyone, we do manage to carry out            No computer ever "wants" to do
   many steps of reasoning, all the             anything,   because   it   was
   higher levels notwithstanding.               programmed by someone else.
   That shows that we humans                    Only if it could program itself
   operate without need of rules: we            from zero on up-an absurdity-
   are "informal systems". On the               would it have its own sense of
   other hand, as an argument against           desire.
   the possibility of any mechanical
   instantiation of reasoning, it is         In his argument, Samuel reconstructs the
   valid,     for  any     mechanical        Tortoise's position, replacing "to reason"
   reasoning-system would have to            by "to want". He implies that behind any
   depend on rules explicitly, and so        mechanization of desire, there has to be
   it couldn't get off the ground            either an infinite regress or worse, a
   unless it had metarules telling it        closed loop. If this is why computers
   when to apply its rules,                  have no will of their own, what about
   metametarules telling it when to          people? The same criterion would imply
   apply its metarules, and so on. We        that
   may conclude that the ability to
   reason can never be mechanized. It
   is a uniquely human capability.

What is wrong with this Devil's advocate
point of view? It is obviously the
assumption that a machine cannot do

Strange Loops, Or Tangled Hierarchies                                              682

  Unless a person designed himself                  Now Samuel's statement brought
  and chose his own wants (as well          up a concept which I want to pursue. It is
  as choosing to choose his own             this: When we humans think, we
  wants, etc.), he cannot be said to        certainly do change our own mental
  have a will of his own.                   rules, and we change the rules that
                                            change the rules, and on and on-but
It makes you pause to think where your      these are, so to speak, "software rules".
sense of having a will comes from.          However, the rules at bottom do not
Unless you are a soulist, you'll probably   change. Neurons run in the same simple
say that it comes from your brain-a piece   way the whole time. You can't "think"
of hardware which you did not design or     your neurons into running some
choose. And yet that doesn't diminish       nonneural way, although you can make
your sense that you want certain things,    your mind change style or subject of
and not others. You aren't a "self-         thought. Like Achilles in the Prelude,
programmed object" (whatever that           Ant Fugue, you have access to your
would be), but you still do have a sense    thoughts but not to your neurons.
of desires, and it springs from the         Software rules on various levels can
physical substrate of your mentality.       change; hardware rules cannot-in fact, to
Likewise, machines may someday have         their rigidity is due the software's
wills despite the fact that no magic        flexibility! Not a paradox at all, but a
program spontaneously appears in            fundamental, simple fact about the
memory from out of nowhere (a "self-        mechanisms of intelligence.
programmed program"). They will have                This distinction between self-
wills for much the same reason as you       modifiable software and inviolate
do-by reason of organization and            hardware is what I wish to pursue in this
structure on many levels of hardware        final Chapter, developing it into a set of
and software. Moral: The Samuel             variations on a theme. Some of the
argument doesn't say anything about the     variations may seem to be quite far-
differences    between     people     and   fetched, but I hope that by the time I
machines, after all. (And indeed, will      close the loop by returning to brains,
will be mechanized.)                        minds,      and     the     sensation   of
                                            consciousness, you will have found an
Below Every Tangled Hierarchy               invariant core in all the variations.
    Lies An Inviolate Level                         My main aim in this Chapter is to
                                            communicate some of the images which
Right after the Two-Part Invention, I       help me to visualize how consciousness
wrote that a central issue of this book     rises out of the jungle of neurons; to
would be: "Do words and thoughts            communicate a set of intangible
follow formal rules?" One major thrust      intuitions, in the hope that
of the book has been to point out the
many-leveledness of the mind/brain, and
I have tried to show why the ultimate
answer to the question is, "Yes-provided
that you go down to the lowest level-the
hardware-to find the rules."


Strange Loops, Or Tangled Hierarchies                                             683

                                               checkerboard. If you can devise a simple
these intuitions are valuable and may          formal notation for expressing rules and
perhaps help others a l4tle to come to         metarules, then to manipulate them will
clearer formulations of their own images       be like manipulating strings formally, or
of what makes minds run. I could not           even like manipulating chess pieces. To
hope for more than that my own mind's          carry things to their logical extreme, you
blurry images of minds and images              could even express rules and metarules
should catalyze the formation of sharper       as positions on auxiliary chess boards.
images of minds and images in other            Then an arbitrary chess position could be
minds.                                         read as a game, or as a set of rules, or as
                                               a set of metarules, etc., depending on
      A Self-Modifying Game                    which interpretation you place on it. Of
                                               course, both players would have to agree
A first variation, then, concerns games in     on conventions for interpreting the
which on your turn, you may modify the         notation.
rules. Think of chess. Clearly the rules               Now we can have any number of
stay the same, just the board position         adjacent chess boards: one for the game,
changes on each move. But let's invent a       one for rules, one for metarules, one for
variation in which, on your turn, you can      metametarules, and so on, as far as you
either make a move or change the rules.        care to carry it. On your turn, you may
But how? At liberty? Can you turn it into      make a move on any one of the chess
checkers? Clearly such anarchy would           boards except the top-level one, using
be pointless. There must be some               the rules which apply (they come from
constraints. For instance, one version         the next chess board up in the hierarchy).
might allow you to redefine the knight's       Undoubtedly both players would get
move. Instead of being 1-and-then-2, it        quite disoriented by the fact that almost
could be m-and-then-n where m and n            anything-though not everything!-can
are arbitrary natural numbers; and on          change. By definition, the top-level
your turn you could change either m or n       chess board can't be changed, because
by plus or minus 1.-So it could go from        you don't have rules telling how to
1-2 to 1-3 to 0-3 to 0-4 to 0-5 to 1-5 to 2-   change it. It is inviolate. There is more
5 ... Then there could be rules about          that is inviolate: the conventions by
redefining the bishop's moves, and the         which the different boards are
other pieces' moves as well. There could       interpreted, the agreement to take turns,
be rules about adding new squares, or          the agreement that each person may
deleting old squares .. .                      change one chess board each turn-and
        Now we have two layers of rules:       you will find more if you examine the
those which tell how to move pieces, and       idea carefully.
those which tell how to change the rules.
So we have rules and metarules. The
next step is obvious: introduce
metametarules by which we can change
the metarules. It is not so obvious how to
do this. The reason it is easy to
formulate rules for moving pieces is that
pieces move in a formalized space: the

Strange Loops, Or Tangled Hierarchies                                                 684

Now it is possible to go considerably                   As you have no doubt imagined,
further in removing the pillars by which        there is nothing to stop us from doing the
orientation is achieved. One step at a          "impossible"-namely, tangling the I-
time. .. We begin by collapsing the             level and the T-level by making the
whole array of boards into a single             interpretation conventions themselves
board. What is meant by this? There will        subject to revision, according to the
be two ways of interpreting the board:          position on the chess board. But in order
(1) as pieces to be moved; (2) as rules         to carry out such a "supertangling",
for moving the pieces. On your turn, you        you'd have to agree with your opponent
move pieces-and perforce, you change            on some further conventions connecting
rules! Thus, the rules constantly change        the two levels-and the act of doing so
themselves. Shades of Typogenetics-or           would create a new level, a new sort of
for that matter, of real genetics. The          inviolate level on top of the
distinction between game, rules,                "supertangled" level (or underneath it, if
metarules, metametarules, has been lost.        you prefer). And this could continue
What was once a nice clean hierarchical         going on and on. In fact, the `jumps"
setup has become a Strange Loop, Or             which are being made are very similar to
Tangled Hierarchy. The moves change             those charted in the Birthday
the rules, the rules determine the moves,       Cantatatata, and in the repeated
round and round the mulberry bush ...           Gödelization      applied     to    various
There are still different levels, but the       improvements on TNT. Each time you
distinction between "lower" and "higher"        think you have reached the end, there is
has been wiped out.                             some new variation on the theme of
        Now, part of what was inviolate         jumping out of the system which
has been made changeable. But there is          requires a kind of creativity to spot.
still plenty that is inviolate. Just as
before, there are conventions between           The Authorship Triangle Again
you and your opponent by which you
interpret the board as a collection of          But I am not interested in pursuing the
rules. There is the agreement to take           strange topic of the ever more abstruse
turns-and probably other implicit               tanglings which can arise in self-
conventions, as well. Notice, therefore,        modifying chess. The point of this has
that the notion of different levels has         been to show, in a somewhat graphic
survived, in an unexpected way. There is        way, how in any system there is always
an Inviolate level-let's call it the I-level-   some "protected" level which is
on which the interpretation conventions         unassailable by the rules on other levels,
reside; there is also a Tangled level-the       no matter how tangled their interaction
T-level-on which the Tangled Hierarchy          may be among themselves. An amusing
resides. So these two levels are still          riddle from Chapter IV illustrates this
hierarchical: the I-level governs what          same idea in a slightly different context.
happens on the T-level, but the T-level         Perhaps it will catch you off guard:
does not and cannot affect the I-level.
No matter that the T-level itself is a
Tangled Hierarchy-it is still governed by
a set of conventions outside of itself.
And that is the important point.

Strange Loops, Or Tangled Hierarchies                                                  685

                                                FIGURE 134. An “authorship triangle”




  There are three authors-Z, T, and E. Now it happens that Z exists only in a novel by
  T. Likewise, T exists only in a novel by E. And strangely, E, too, exists only in a
  novel-by Z, of course. Now, is such an "authorship triangle" really possible? (See
  Fig. 134.)

Of course it's possible. But there's a trick ... All three authors Z, T, E, are themselves
characters in another novel-by H! You can think of the Z-T-E triangle as a Strange Loop,
Or Tangled Hierarchy; but author H is outside of the space in which that tangle takes
place-author H is in an inviolate space. Although Z, T, and E all have access-direct or
indirect-to each other, and can do dastardly things to each other in their various novels,
none of them can touch H's life! They can't even imagine him-no more than you can
imagine the author of the book you're a character in. If I were to draw author H, I would
represent him somewhere off the page. Of course that would present a problem, since
drawing a thing necessarily puts it onto the page ... Anyway, H is really outside of the
world of Z, T, and E, and should be represented as being so.

                             Escher's Drawing Hands
Another classic variation on our theme is the Escher picture of Drawing Hands (Fig.
135). Here, a left hand (LH) draws a right hand (RH), while at the same time, RH draws
LH. Once again, levels which ordinarily are seen as hierarchical-that which draws, and
that which is drawn-turn back on each other, creating a Tangled Hierarchy. But the theme
of the Chapter is borne out, of course, since behind it all lurks the undrawn but drawing
hand of M. C. Escher, creator of both LH and RH. Escher is outside of the two-hand
space, and in my schematic version of his picture (Fig. 136), you can see that explicitly.
In this schematized representation of the Escher picture, you see the Strange Loop, Or
Tangled Hierarchy at the top; also, you see the Inviolate Level below it, enabling it to
come into being. One could further Escherize the Escher picture, by taking a photograph
of a hand drawing it. And so on.




Strange Loops, Or Tangled Hierarchies                                                 686

          FIGURE 135. Drawing Hands, by M. C. Escher (lithograph, 1948).




 FIGURE 136. Abstract diagram of M. C. Escher's Drawing Hands. On top, a seeming
                          paradox. Below, its resolution.

Strange Loops, Or Tangled Hierarchies                                         685

                              Brain and Mind:
                 A Neural Tangle Supporting a Symbol Tangle

Now we can relate this to the brain, as well as to Al programs. In our thoughts, symbols
activate other symbols, and all interact heterarchically. Furthermore, the symbols may
cause each other to change internally, in the fashion of programs acting on other
programs. The illusion is created, because of the Tangled Hierarchy of symbols, that
there is no inviolate level. one thinks there is no such level because that level is shielded
from our view.
        If it were possible to schematize this whole image, there would be a gigantic
forest of symbols linked to each other by tangly lines like vines in a tropical jungle-this
would be the top level, the Tangled Hierarchy where thoughts really flow back and forth.
This is the elusive level of mind: the analogue to LH and RH. Far below in the schematic
picture, analogous to the invisible "prime mover" Escher, there would be a representation
of the myriad neurons-the "inviolate substrate" which lets the tangle above it come into
being. Interestingly, this other level is itself a tangle in a literal sense-billions of cells and
hundreds of billions of axons, joining them all together.
        This is an interesting case where a software tangle, that of the symbols, is
supported by a hardware tangle, that of the neurons. But only the symbol tangle is a
Tangled Hierarchy. The neural tangle is* ust a "simple" tangle. This distinction is pretty
much the same as that between Strange Loops and feedback, which I mentioned in
Chapter XVI. A Tangled Hierarchy occurs when what you presume are clean hierarchical
levels take you by surprise and fold back in a hierarchy-violating way. The surprise
element is important; it is the reason I call Strange Loops "strange". A simple tangle, like
feedback, doesn't involve violations of presumed level distinctions. An example is when
you're in the shower and you wash your left arm with your right, and then vice versa.
There is no strangeness to the image. Escher didn't choose to draw hands drawing hands
for nothing!
        Events such as two arms washing each other happen all the time in the world, and
we don't notice them particularly. I say something to you, then you say something back to
me. Paradox \% No; our perceptions of each other didn't involve a hierarchy to begin with,
so there is no sense of strangeness.
        On the other hand, where language does create strange loops is when it talks
about itself, whether directly or indirectly. Here, something in the system jumps out and
acts on the system, as if it were outside the system. What bothers us is perhaps an ill-
defined sense of topological wrongness: the inside-outside distinction is being blurred, as
in the famous shape called a "Klein bottle". Even though the system is an abstraction, our
minds use spatial imagery with a sort of mental topology.
        Getting back to the symbol tangle, if we look only at it, and forget the neural
tangle, then we seem to see a self-programmed object-in just the same way as we seem to
see a self-drawn picture if we look at Drawing Hands and somehow fall for the illusion,
by forgetting the existence of Escher. For




Strange Loops, Or Tangled Hierarchies                                                        686

the picture, this is unlikely-but for humans and the way they look at their minds, this is
usually what happens. We feel self-programmed. Indeed, we couldn't feel any other way,
for we are shielded from the lower levels, the neural tangle. Our thoughts seem to run
about in their own space, creating new thoughts and modifying old ones, and we never
notice any neurons helping us out! But that is to be expected. We can't.
        An analogous double-entendre can happen with LISP programs that are designed
to reach in and change their own structure. If you look at them on the LISP level, you
will say that they change themselves; but if you shift levels, and think of LISP programs
as data to the LISP interpreter (see Chapter X), then in fact the sole program that is
running is the interpreter, and the changes being made are merely changes in pieces of
data. The LISP interpreter itself is shielded from changes.
        How you describe a tangled situation of this sort depends how far back you step
before describing. If you step far enough back, you can often see the clue that allows you
to untangle things.

                          Strange Loops in Government
A fascinating area where hierarchies tangle is government-particularly in the courts.
Ordinarily, you think of two disputants arguing their cases in court, and the court
adjudicating the matter. The court is on a different level from the disputants. But strange
things can start to happen when the courts themselves get entangled in legal cases.
Usually there is a higher court which is outside the dispute. Even if two lower courts get
involved in some sort of strange fight, with each one claiming jurisdiction over the other,
some higher court is outside, and in some sense it is analogous to the inviolate
interpretation conventions which we discussed in the warped version of chess.
         But what happens when there is no higher court, and the Supreme Court itself gets
all tangled up in legal troubles? This sort of snarl nearly happened in the Watergate era.
The President threatened to obey only a "definitive ruling" of the Supreme Court-then
claimed he had the right to decide what is "definitive". Now that threat never was made
good; but if it had been, it would have touched off a monumental confrontation between
two levels of government, each of which, in some ways, can validly claim to be "above"
the other-and to whom is there recourse to decide which one is right? To say "Congress"
is not to settle the matter, for Congress might command the President to obey the
Supreme Court, yet the President might still refuse, claiming that he has the legal right to
disobey the Supreme Court (and Congress!) under certain circumstances. This would
create a new court case, and would throw the whole system into disarray, because it
would be so unexpected, so Tangled-so Strange!
         The irony is that once you hit your head against the ceiling like this, where you
are prevented from jumping out of the system to a yet higher authority, the only recourse
is to forces which seem less well defined by




Strange Loops, Or Tangled Hierarchies                                                   687

rules, but which are the only source of higher-level rules anyway: the lower-level rules,
which in this case means the general reaction of society. It is well to remember that in a
society like ours, the legal system is, in a sense, a polite gesture granted collectively by
millions of people-and it can be overridden just as easily as a river can overflow its
banks. Then a seeming anarchy takes over; but anarchy has its own kinds of rules, no less
than does civilized society: it is just that they operate from the bottom up, not from the
top down. A student of anarchy could try to discover rules according to which anarchic
situations develop in time, and very likely there are some such rules.
         An analogy from physics is useful here. As was mentioned earlier in the book,
gases in equilibrium obey simple laws connecting their temperature, pressure, and
volume. However, a gas can violate those laws (as a President can violate laws)-provided
it is not in a state of equilibrium. In nonequilibrium situations, to describe what happens,
a physicist has recourse only to statistical mechanics-that is, to a level of description
which is not macroscopic, for the ultimate explanation of a gas's behavior always lies on
the molecular level, just as the ultimate explanation of a society's political behavior
always lies at the "grass roots level". The field of nonequilibrium thermodynamics
attempts to find macroscopic laws to describe the behavior of gases (and other systems)
which are out of equilibrium. It is the analogue to the branch of political science which
would search for laws governing anarchical societies.
         Other curious tangles which arise in government include the FBI investigating its
own wrongdoings, a sheriff going to jail while in office, the self-application of the
parliamentary rules of procedure, and so on. One of the most curious legal cases I ever
heard of involved a person who claimed to have psychic powers. In fact, he claimed to be
able to use his psychic powers to detect personality traits, and thereby to aid lawyers in
picking juries. Now what if this "psychic" has to stand trial himself one day? What effect
might this have on a jury member who believes staunchly in ESP? How much will he feel
affected by the psychic (whether or not the psychic is genuine)? The territory is ripe for
exploitation-a great area for selffulfilling prophecies.

                   Tangles Involving Science and the Occult
Speaking of psychics and ESP, another sphere of life where strange loops abound is
fringe science. What fringe science does is to call into question many of the standard
procedures or beliefs of orthodox science, and thereby challenge the objectivity of
science. New ways of interpreting evidence that rival the established ones are presented.
But how do you evaluate a way of interpreting evidence? Isn't this precisely the problem
of objectivity all over again, just on a higher plane? Of course. Lewis Carroll's infinite-
regress paradox appears in a new guise. The Tortoise would argue that if you want to
show that A is a fact, you need evidence: B. But what makes you sure that B is evidence
of A?' To show that, you need meta-




Strange Loops, Or Tangled Hierarchies                                                   688

evidence: C. And for the validity of that meta-evidence, you need metameta-evidence-
and so on, ad nauseam. Despite this argument, people have an intuitive sense of evidence.
This is because-to repeat an old refrain-people have built-in hardware in their brains that
includes some rudimentary ways of interpreting evidence. We can build on this, and
accumulate new ways of interpreting evidence; we even learn how and when to override
our most basic mechanisms of evidence interpretation, as one must, for example, in trying
to figure out magic tricks.
        Concrete examples of evidence dilemmas crop up in regard to many phenomena
of fringe science. For instance, ESP often seems to manifest itself outside of the
laboratory, but when brought into the laboratory, it vanishes mysteriously. The standard
scientific explanation for this is that ESP is a nonreal phenomenon which cannot stand up
to rigorous scrutiny. Some (by no means all) believers in ESP have a peculiar way of
fighting back, however. They say, "No, ESP is real; it simply goes away when one tries
to observe it scientifically-it is contrary to the nature of a scientific worldview." This is
an amazingly brazen technique, which we might call "kicking the problem upstairs".
What that means is, instead of questioning the matter at hand, you call into doubt theories
belonging to a higher level of credibility. The believers in ESP insinuate that what is
wrong is not their ideas, but the belief system of science. This is a pretty grandiose claim,
and unless there is overwhelming evidence for it, one should be skeptical of it. But then
here we are again, talking about "overwhelming evidence" as if everyone agreed on what
that means!

                               The Nature of Evidence
The Sagredo-Simplicio-Salviati tangle, mentioned in Chapters XIII and XV, gives
another example of the complexities of evaluation of evidence. Sagredo tries to find some
objective compromise, if possible, between the opposing views of Simplicio and Salviati.
But compromise may not always be possible. How can one compromise "fairly" between
right and wrong? Between fair and unfair? Between compromise and no compromise?
These questions come up over and over again in disguised form in arguments about
ordinary things.
        Is it possible to define what evidence is? Is it possible to lay down laws as to how
to make sense out of situations? Probably not, for any rigid rules would undoubtedly have
exceptions, and nonrigid rules are not rules. Having an intelligent AI program would not
solve the problem either, for as an evidence processor, it would not be any less fallible
than humans are. So, if evidence is such an intangible thing after all, why am I warning
against new ways of interpreting evidence? Am I being inconsistent? In this case, I don't
think so. My feeling is that there are guidelines which one can give, and out of them an
organic synthesis can be made. But inevitably some amount of judgment and intuition
must enter the picture-things which are different in different people. They will also be
different in




Strange Loops, Or Tangled Hierarchies                                                    689

different AI programs. Ultimately, there are complicated criteria for deciding if a method
of evaluation of evidence is good. One involves the "usefulness" of ideas which are
arrived at by that kind of reasoning. Modes of thought which lead to useful new things in
life are deemed "valid" in some sense. But this word "useful" is extremely subjective.
         My feeling is that the process by which we decide what is valid or what is true is
an art; and that it relies as deeply on a sense of beauty and simplicity as it does on rock-
solid principles of logic or reasoning or anything else which can be objectively
formalized. I am not saying either (1) truth is a chimera, or (2) human intelligence is in
principle not programmable. I am saying (1) truth is too elusive for any human or any
collection of humans ever to attain fully; and (2) Artificial Intelligence, when it reaches
the level of human intelligence-or even if it surpasses it-will still be plagued by the
problems of art, beauty, and simplicity, and will run up against these things constantly in
its own search for knowledge and understanding.
         "What is evidence?" is not just a philosophical question, for it intrudes into life in
all sorts of places. You are faced with an extraordinary number of choices as to how to
interpret evidence at every moment. You can hardly go into a bookstore (or these days,
even a grocery store!) without seeing books on clairvoyance, ESP, UFO's, the Bermuda
triangle, astrology, dowsing, evolution versus creation, black holes, psi fields,
biofeedback, transcendental meditation, new theories of psychology ... In science, there
are fierce debates about catastrophe theory, elementary particle theory, black holes, truth
and existence in mathematics, free will, Artificial Intelligence, reductionism versus
holism ... On the more pragmatic side of life, there are debates over the efficacy of
vitamin C or of laetrile, over the real size of oil reserves (either underground or stored),
over what causes inflation and unemployment-and on and on. There is Buckminster
Fullerism, Zen Buddhism, Zeno's paradoxes, psychoanalysis, etc., etc. From issues as
trivial as where books ought to be shelved in a store, to issues as vital as what ideas are to
be taught to children in schools, ways of interpreting evidence play an inestimable role.

                                     Seeing Oneself
One of the most severe of all problems of evidence interpretation is that of trying to
interpret all the confusing signals from the outside as to who one is. In this case, the
potential for intralevel and interlevel conflict is tremendous. The psychic mechanisms
have to deal simultaneously with the individual's internal need for self-esteem and the
constant flow of evidence from the outside affecting the self-image. The result is that
information flows in a complex swirl between different levels of the personality; as it
goes round and round, parts of it get magnified, reduced, negated, or otherwise distorted,
and then those parts in turn get further subjected to the same sort of swirl, over and over
again-all of this in an attempt to reconcile what is, with what we wish were (see Fig. 81).




Strange Loops, Or Tangled Hierarchies                                                      690

        The upshot is that the total picture of "who I am" is integrated in some
enormously complex way inside the entire mental structure, and contains in each one of
us a large number of unresolved, possibly unresolvable, inconsistencies. These
undoubtedly provide much of the dynamic tension which is so much a part of being
human. Out of this tension between the inside and outside notions of who we are come
the drives towards various goals that make each of us unique. Thus, ironically, something
which we all have in common-the fact of being self-reflecting conscious beings-leads to
the rich diversity in the ways we have of internalizing evidence about all sorts of things,
and in the end winds up being one of the major forces in creating distinct individuals.

                     Gödel’s Theorem and Other Disciplines
It is natural to try to draw parallels between people and sufficiently complicated formal
systems which, like people, have "self-images" of a sort. Gödel’s Theorem shows that
there are fundamental limitations to consistent formal systems with self-images. But is it
more general? Is there a "Gödel’s Theorem of psychology", for instance?
        If one uses Gödel’s Theorem as a metaphor, as a source of inspiration, rather than
trying to translate it literally into the language of psychology or of any other discipline,
then perhaps it can suggest new truths in psychology or other areas. But it is quite
unjustifiable to translate it directly into a statement of another discipline and take that as
equally valid. It would be a large mistake to think that what has been worked out with the
utmost delicacy in mathematical logic should hold without modification in a completely
different area.

               Introspection and Insanity: A Gödelian Problem
I think it can have suggestive value to translate Gödel’s Theorem into other domains,
provided one specifies in advance that the translations are metaphorical and are not
intended to be taken literally. That having been said, I see two major ways of using
analogies to connect Gödel’s Theorem and human thoughts. One involves the problem of
wondering about one's sanity. How can you figure out if you are sane? This is a Strange
Loop indeed. Once you begin to question your own sanity, you can get trapped in an
ever-tighter vortex of self-fulfilling prophecies, though the process is by no means
inevitable. Everyone knows that the insane interpret the world via their own peculiarly
consistent logic; how can you tell if your own logic is "peculiar" or not, given that you
have only your own logic to judge itself? I don't see any answer. I am just reminded of
Gödel’s second Theorem, which implies that the only versions of formal number theory
which assert their own consistency are inconsistent ...




Strange Loops, Or Tangled Hierarchies                                                     691

              Can We Understand Our Own" Minds or Brains?
The other metaphorical analogue to Gödel’s Theorem which I find provocative suggests
that ultimately, we cannot understand our own minds/ brains. This is such a loaded,
many-leveled idea that one must be extremely cautious in proposing it. What does
"understanding our own minds/brains" mean? It could mean having a general sense of
how they work, as mechanics have a sense of how cars work. It could mean having a
complete explanation for why people do any and all things they do. It could mean having
a complete understanding of the physical structure of one's own brain on all levels. It
could mean having a complete wiring diagram of a brain in a book (or library or
computer). It could mean knowing, at every instant, precisely what is happening in one's
own brain on the neural level-each firing, each synaptic alteration, and so on. It could
mean having written a program which passes the Turing test. It could mean knowing
oneself so perfectly that such notions as the subconscious and the intuition make no
sense, because everything is out in the open. It could mean any number of other things.
        Which of these types of self-mirroring, if any, does the self-mirroring in Gödel’s
Theorem most resemble? I would hesitate to say. Some of them are quite silly. For
instance, the idea of being able to monitor your own brain state in all its detail is a pipe
dream, an absurd and uninteresting proposition to start with; and if Gödel’s Theorem
suggests that it is impossible, that is hardly a revelation. On the other hand, the age-old
goal of knowing yourself in some profound way-let us call it "understanding your own
psychic structure"-has a ring of plausibility to it. But might there not be some vaguely
Godelian loop which limits the depth to which any individual can penetrate into his own
psyche? Just as we cannot see our faces with our own eyes, is it not reasonable to expect
that we cannot mirror our complete mental structures in the symbols which carry them
out?
        All the limitative Theorems of metamathematics and the theory of computation
suggest that once the ability to represent your own structure has reached a certain critical
point, that is the kiss of death: it guarantees that you can never represent yourself totally.
Gödel’s Incompleteness Theorem, Church's Undecidability Theorem, Turing's Halting
Theorem, Tarski's Truth Theorem-all have the flavor of some ancient fairy tale which
warns you that "To seek self-knowledge is to embark on a journey which ... will always
be incomplete, cannot be charted on any map, will never halt, cannot be described."
But do the limitative Theorems have any bearing on people? Here is one way of arguing
the case. Either I am consistent or I am inconsistent. (The latter is much more likely, but
for completeness' sake, I consider both possibilities.) If I am consistent, then there are
two cases. (1) The "low-fidelity" case: my self-understanding is below a certain critical
point. In this case, I am incomplete by hypothesis. (2) The "high-fidelity" case: My self-
understanding has reached the critical point where a metaphorical analogue of the
limitative Theorems does apply, so my self-understanding




Strange Loops, Or Tangled Hierarchies                                                     692

undermines itself in a Gödelian way, and I am incomplete for that reason. Cases (1) and
(2) are predicated on my being 100 per cent consistent-a very unlikely state of affairs.
More likely is that I am inconsistent-but that's worse, for then inside me there are
contradictions, and how can I ever understand that?
        Consistent or inconsistent, no one is exempt from the mystery of the self.
Probably we are all inconsistent. The world is just too complicated for a person to be able
to afford the luxury of reconciling all of his beliefs with each other. Tension and
confusion are important in a world where many decisions must be made quickly, Miguel
de Unamuno once said, "If a person never contradicts himself, it must be that he says
nothing." I would say that we all are in the same boat as the Zen master who, after
contradicting himself several times in a row, said to the confused Doko, "I cannot
understand myself."

                 Gödel’s Theorem and Personal Nonexistence
Perhaps the greatest contradiction in our lives, the hardest to handle, is the knowledge
"There was a time when I was not alive, and there will come a time when I am not alive."
On one level, when you "step out of yourself" and see yourself as "just another human
being", it makes complete sense. But on another level, perhaps a deeper level, personal
nonexistence makes no sense at all. All that we know is embedded inside our minds, and
for all that to be absent from the universe is not comprehensible. This is a basic
undeniable problem of life; perhaps it is the best metaphorical analogue of Gödel’s
Theorem. When you try to imagine your own nonexistence, you have to try to jump out
of yourself, by mapping yourself onto someone else. You fool yourself into believing that
you can import an outsider's view of yourself into you, much as TNT "believes" it
mirrors its own metatheory inside itself. But TNT only contains its own metatheory up to
a certain extent-not fully. And as for you, though you may imagine that you have jumped
out of yourself, you never can actually do so-no more than Escher's dragon can jump out
of its native two-dimensional plane into three dimensions. In any case, this contradiction
is so great that most of our lives we just sweep the whole mess under the rug, because
trying to deal with it just leads nowhere.
        Zen minds, on the other hand, revel in this irreconcilability. Over and over again,
they face the conflict between the Eastern belief: "The world and I are one, so the notion
of my ceasing to exist is a contradiction in terms" (my verbalization is undoubtedly too
Westernized-apologies to Zenists), and the Western belief: "I am just part of the world,
and I will die, but the world will go on without me."

                                Science and Dualism
Science is often criticized as being too "Western" or "dualistic"-that is, being permeated
by the dichotomy between subject and object, or observer




Strange Loops, Or Tangled Hierarchies                                                  693

and observed. While it is true that up until this century, science was exclusively
concerned with things which can be readily distinguished from their human observers-
such as oxygen and carbon, light and heat, stars and planets, accelerations and orbits, and
so on-this phase of science was a necessary prelude to the more modern phase, in which
life itself has come under investigation. Step by step, inexorably, "Western" science has
moved towards investigation of the human mind-which is to say, of the observer.
Artificial Intelligence research is the furthest step so far along that route. Before AI came
along, there were two major previews of the strange consequences of the mixing of
subject and object in science. One was the revolution of quantum mechanics, with its
epistemological problems involving the interference of the observer with the observed.
The other was the mixing of subject and object in metamathematics, beginning with
Godel's Theorem and moving through all the other limitative'Theorems we have
discussed. Perhaps the next step after Al will be the self-application of science: science
studying itself as an object. This is a different manner of mixing subject and object-
perhaps an even more tangled one than that of humans studying their own minds.
         By the way, in passing, it is interesting to note that all results essentially
dependent on the fusion of subject and object have been limitative results. In addition to
the limitative Theorems, there is Heisenberg's uncertainty principle, which says that
measuring one quantity renders impossible the simultaneous measurement of a related
quantity. I don't know why all these results are limitative. Make of it what you will.

                 Symbol vs. Object in Modern Music and Art
Closely linked with the subject-object dichotomy is the symbol-object dichotomy, which
was explored in depth by Ludwig Wittgenstein in the early part of this century. Later the
words "use" and "mention" were adopted to make the same distinction. Quine and others
have written at length about the connection between signs and what they stand for. But
not only philosophers have devoted much thought to this deep and abstract matter. In our
century both music and art have gone through crises which reflect a profound concern
with this problem. Whereas music and painting, for instance, have traditionally expressed
ideas or emotions through a vocabulary of "symbols" (i.e. visual images, chords,
rhythms, or whatever), now there is a tendency to explore the capacity of music and art to
not express anything just to be. This means to exist as pure globs of paint, or pure sounds,
but in either case drained of all symbolic value.
        In music, in particular, John Cage has been very influential in bringing a Zen-like
approach to sound. Many of his pieces convey a disdain for "use" of sounds-that is, using
sounds to convey emotional states-and an exultation in "mentioning" sounds-that is,
concocting arbitrary juxtapositions of sounds without regard to any previously formulated
code by which a listener could decode them into a message. A typical example is
"Imaginary Landscape no. 4", the polyradio piece described in Chapter VI. I may not




Strange Loops, Or Tangled Hierarchies                                                    694

be doing Cage justice, but to me it seems that much of his work has been directed at
bringing meaninglessness into music, and in some sense, at making that meaninglessness
have meaning. Aleatoric music is a typical exploration in that direction. (Incidentally,
chance music is a close cousin to the much later notion of "happenings" or "be-in"' s.)
There are many other contemporary composers who are following Cage's lead, but few
with as much originality. A piece by Anna Lockwood, called "Piano Burning", involves
just that-with the strings stretched to maximum tightness, to make them snap as loudly as
possible; in a piece by LaMonte Young, the noises are provided by shoving the piano all
around the stage and through obstacles, like a battering ram.
        Art in this century has gone through many convulsions of this general type. At
first there was the abandonment of representation, which was genuinely revolutionary:
the beginnings of abstract art. A gradual swoop from pure representation to the most
highly abstract patterns is revealed in the work of Piet Mondrian. After the world was
used to nonrepresentational art, then surrealism came along. It was a bizarre about-face,
something like neoclassicism in music, in which extremely representational art was
"subverted" and used for altogether new reasons: to shock, confuse, and amaze. This
school was founded by Andre Breton, and was located primarily in France; some of its
more influential members were Dali, Magritte, de Chirico, Tanguy.

                          Magritte's Semantic Illusions
Of all these artists, Magritte was the most conscious of the symbol-object mystery (which
I see as a deep extension of the use-mention distinction). He uses it to evoke powerful
responses in viewers, even if the viewers do not verbalize the distinction this way. For
example, consider his very strange variation on the theme of still life, called Common
Sense (Fig. 137).

              FIGURE 137. Common Sense, by Rene Magritte (1945-46).




Strange Loops, Or Tangled Hierarchies                                                695

                FIGURE 138. The Two Mysteries, by Rene Magritte (1966).

Here, a dish filled with fruit, ordinarily the kind of thing represented inside a still life, is
shown sitting on top of a blank canvas. The conflict between the symbol and the real is
great. But that is not the full irony, for of course the whole thing is itself just a painting-in
fact, a still life with nonstandard subject matter.
         Magritte's series of pipe paintings is fascinating and perplexing. Consider The
Two Mysteries (Fig. 138). Focusing on the inner painting, you get the message that
symbols and pipes are different. Then your glance moves upward to the "real" pipe
floating in the air-you perceive that it is real, while the other one is just a symbol. But that
is of course totally wrong: both of them are on the same flat surface before your eyes.
The idea that one pipe is in a twice-nested painting, and therefore somehow "less real"
than the other pipe, is a complete fallacy. Once you are willing to "enter the room", you
have already been tricked: you've fallen for image as reality. To be consistent in your
gullibility, you should happily go one level further down, and confuse image-within-
image with reality. The only way not to be sucked in is to see both pipes merely as
colored smudges on a surface a few inches in front of your nose. Then, and only then, do
you appreciate the full meaning of the written message "Ceci West pas une pipe"-but
ironically, at the very instant everything turns to smudges, the writing too turns to
smudges, thereby losing its meaning! In other words, at that instant, the verbal message
of the painting self-destructs in a most Gödelian way.




Strange Loops, Or Tangled Hierarchies                                                       696

                  FIGURE 139. Smoke Signal. [Drawing by the author.]

         The Air and the Song (Fig. 82), taken from a series by Magritte, accomplishes all
that The Two Mysteries does, but in one level instead of two. My drawings Smoke Signal
and Pipe Dream (Figs. 139 and 140) constitute "Variations on a Theme of Magritte". Try
staring at Smoke Signal for a while. Before long, you should be able to make out a hidden
message saying, "Ceci n'est pas un message". Thus, if you find the message, it denies
itself-yet if you don't, you miss the point entirely. Because of their indirect self-snuffing,
my two pipe pictures can be loosely mapped onto Gödel’s G-thus giving rise to a
"Central Pipemap", in the same spirit as the other "Central Xmaps": Dog, Crab, Sloth.
         A classic example of use-mention confusion in paintings is the occurrence of a
palette in a painting. Whereas the palette is an illusion created by the representational
skill of the painter, the paints on the painted palette are literal daubs of paint from the
artist's palette. The paint plays itself-it does not symbolize anything else. In Don
Giovanni, Mozart exploited a related trick: he wrote into the score explicitly the sound of
an orchestra tuning up. Similarly, if I want the letter 'I' to play itself (and not symbolize
me), I put 'I' directly into my text; then I enclose `I' between quotes. What results is ''I"
(not `I', nor "`I"'). Got that?




Strange Loops, Or Tangled Hierarchies                                                     697

                  FIGURE 140. Pipe Dream. [Drawing by the author.]

                            The "Code" of Modern Art
A large number of influences, which no one could hope to pin down completely, led to
further explorations of the symbol-object dualism in art. There is no doubt that John
Cage, with his interest in Zen, had a profound influence on art as well as on music. His
friends jasper Johns and Robert Rauschenberg both explored the distinction between
objects and symbols by using objects as symbols for themselves-or, to flip the coin, by
using symbols as objects in themselves. All of this was perhaps intended to break down
the notion that art is one step removed from reality-that art speaks in "code", for which
the viewer must act as interpreter. The idea was to eliminate the step of interpretation and
let the naked object simply be, period. ("Period"-a curious case of use-mention blur.)
However, if this was the intention, it was a monumental flop, and perhaps had to be.
        Any time an object is exhibited in a gallery or dubbed a "work", it acquires an
aura of deep inner significance-no matter how much the viewer has been warned not to
look for meaning. In fact, there is a backfiring effect whereby the more that viewers are
told to look at these objects without mystification, the more mystified the viewers get.
After all, if a




Strange Loops, Or Tangled Hierarchies                                                   698

wooden crate on a museum floor is just a wooden crate on a museum floor, then why
doesn't the janitor haul it out back and throw it in the garbage? Why is the name of an
artist attached to it? Why did the artist want to demystify art? Why isn't that dirt clod out
front labeled with an artist's name? Is this a hoax? Am I crazy, or are artists crazy? More
and more questions flood into the viewer's mind; he can't help it. This is the "frame
effect" which art-Art-automatically creates. There is no way to suppress the wonderings
in the minds of the curious.
         Of course, if the purpose is to instill a Zen-like sense of the world as devoid of
categories and meanings, then perhaps such art is merely intended to serve-as does
intellectualizing about Zen-as a catalyst to inspire the viewer to go out and become
acquainted with the philosophy which rejects "inner meanings" and embraces the world
as a whole. In this case, the art is self-defeating in the short run, since the viewers do
ponder about its meaning, but it achieves its aim with a few people in the long run, by
introducing them to its sources. But in either case, it is not true that there is no code by
which ideas are conveyed to the viewer. Actually, the code is a much more complex
thing, involving statements about the absence of codes and so forth-that is, it is part code,
part metacode, and so on. There is a Tangled Hierarchy of messages being transmitted by
the most Zen-like art objects, which is perhaps why so many find modern art so
inscrutable.

                                    Ism Once Again
Cage has led a movement to break the boundaries between art and nature. In music, the
theme is that all sounds are equal-a sort of acoustical democracy. Thus silence is just as
important as sound, and random sound is just as important as organized sound. Leonard
B. Meyer, in his book Music, the Arts, and Ideas, has called this movement in music
"transcendentalism", and states:

   If the distinction between art and nature is mistaken, aesthetic valuation is
   irrelevant. One should no more judge the value of a piano sonata than one should
   judge the value of a stone, a thunderstorm, or a starfish. "Categorical statements,
   such as right and wrong, beautiful or ugly, typical of the rationalistic thinking of
   tonal aesthetics," writes Luciano Berio [a contemporary composer, "are no longer
   useful in understanding why and how a composer today works on audible forms
   and musical action."

Later, Meyer continues in describing the philosophical position of transcendentalism:

           ... all things in all of time and space are inextricably connected with one
   another. Any divisions, classifications, or organizations discovered in the universe
   are arbitrary. The world is a complex, continuous, single event .2 [Shades of Zeno!]

       I find "transcendentalism" too bulky a name for this movement. In its place, I use
"ism". Being a suffix without a prefix, it suggests an ideology




Strange Loops, Or Tangled Hierarchies                                                    699

           FIGURE 141. The Human Condition I, by Rene Magritte (1933).




Strange Loops, Or Tangled Hierarchies                                    700

without ideas-which, however you interpret it, is probably the case. And since."ism"
embraces whatever is, its name is quite fitting. In "ism" thL- word "is" is half mentioned,
half used; what could be more appropriate? Ism is the spirit of Zen in art. And just as the
central problem of Zen is to unmask the self, the central problem of art in this century
seems to be to figure out what art is. All these thrashings-about are part of its identity
crisis.
        We have seen that the use-mention dichotomy, when pushed, turns into the
philosophical problem of symbol-object dualism, which links it to the mystery of mind.
Magritte wrote about his painting The Human Condition I (Fig. 141):

   I placed in front of a window, seen from a room, a painting representing exactly
   that part of the landscape which was hidden from view by the painting. Therefore,
   the tree represented in the painting hid from view the tree situated behind it, outside
   the room. It existed for the spectator, as it were, simultaneously in his mind, as both
   inside the room in the painting, and outside in the real landscape. Which is how we
   see the world: we see it as being outside ourselves even though it is only a mental
   representation of it that we
   experience inside ourselves.'

                              Understanding the Mind
First through the pregnant images of his painting, and then in direct words, Magritte
expresses the link between the two questions "How do symbols work?" and "How do our
minds work?" And so he leads us back to the question posed earlier: "Can we ever hope
to understand our minds! brains?"
        Or does some marvelous diabolical Gödelian proposition preclude our ever
unraveling our minds? Provided you do not adopt a totally unreasonable definition of
"understanding", I see no Gödelian obstacle in the way of the eventual understanding of
our minds. For instance, it seems to me quite reasonable to desire to understand the
working principles of brains in general, much the same way as we understand the
working principles of car engines in general. It is quite different from trying to
understand any single brain in every last detail-let alone trying to do this for one's own
brain! I don't see how Gödel’s Theorem, even if construed in the sloppiest way, has
anything to say about the feasibility of this prospect. I see no reason that Gödel’s
Theorem imposes any limitations on our ability to formulate and verify the general
mechanisms by which thought processes take place in the medium of nerve cells. I see no
barrier imposed by Gödel’s Theorem to the implementation on computers (or their
successors) of types of symbol manipulation that achieve roughly the same results as
brains do. It is entirely another question to try and duplicate in a program some particular
human's mind-but to produce an intelligent program at all is a more limited goal. Godel's
Theorem doesn't ban our reproducing our own level of intelligence via programs any
more than it bans our reproducing our own level of intelligence via transmission of
hereditary information in




Strange Loops, Or Tangled Hierarchies                                                    701

DNA, followed by education. Indeed, we have seen, in Chapter XVI, how a remarkable
'Gödelian mechanism-the Strange Loop of proteins and DNA-is precisely what allows
transmission of intelligence!
        Does Gödel’s Theorem, then, have absolutely nothing to offer us in thinking
about our own minds? I think it does, although not in the mystical and [imitative way
which some people think it ought to. I think that the process of coming to understand
Gödel’s proof, with its construction involving arbitrary codes, complex isomorphisms,
high and low levels of interpretation, and the capacity for self-mirroring, may inject some
rich undercurrents and flavors into one's set of images about symbols and symbol
processing, which may deepen one's intuition for the relationship. between mental
structures on different levels.

                       Accidental Inexplicability of Intelligence?

Before suggesting a philosophically intriguing "application" of Godel's proof. I would
like to bring up the idea of "accidental inexplicability" of intelligence. Here is what that
involves. It could be that our brains, unlike car engines, are stubborn and intractable
systems which we cannot neatly decompose in any way. At present, we have no idea
whether our brains will yield to repeated attempts to cleave them into clean layers, each
of which can be explained in terms of lower layers-or whether our brains will foil all our
attempts at decomposition.
        But even if we do fail to understand ourselves, there need not be any Godelian
"twist" behind it; it could be simply an accident of fate that our brains are too weak to
understand themselves. Think of the lowly giraffe, for instance, whose brain is obviously
far below the level required for self-understanding-yet it is remarkably similar to our own
brain. In fact, the brains of giraffes, elephants, baboons-even the brains of tortoises or
unknown beings who are far smarter than we are-probably all operate on basically the
same set of principles. Giraffes may lie far below the threshold of intelligence necessary
to understand how those principles fit together to produce the qualities of mind; humans
may lie closer to that threshold perhaps just barely below it, perhaps even above it. The
point is that there may be no fundamental (i.e., Gödelian) reason why those qualities are
incomprehensible; they may be completely clear to more intelligent beings.

        Undecidability Is Inseparable from a High-Level Viewpoint
Barring this pessimistic notion of the accidental inexplicability of the brain, what insights
might Gödel’s proof offer us about explanations of our minds/brains? Gödel’s proof
offers the notion that a high-level view of a system may contain explanatory power which
simply is absent on the lower levels. By this I mean the following. Suppose someone
gave you G, Gödel’s undecidable string, as a string of TNT. Also suppose you knew
nothing of' Gödel-numbering. The question you are supposed to answer is: "Why isn't




Strange Loops, Or Tangled Hierarchies                                                    702

this string a theorem of TNT?" Now you are used to such questions; for instance, if you
had been asked that question about SO=0, you would have a ready explanation: "Its
negation, ~S0=0, is a theorem." This, together with your knowledge that TNT is
consistent, provides an explanation of why the given string is a nontheorem. This is what
I call an explanation "on the TNT-level". Notice how different it is from the explanation
of why MU is not a theorem of the MIU-system: the former comes from the M-mode, the
latter only from the I-mode.
         Now what about G? The TNT-level explanation which worked for 50=0 does not
work for G, because - G is not a theorem. The person who has no overview of TNT will
be baffled as to why he can't make G according to the rules, because as an arithmetical
proposition, it apparently has nothing wrong with it. In fact, when G is turned into a
universally quantified string, every instance gotten from G by substituting numerals for
the variables can be derived. The only way to explain G's nontheoremhood is to discover
the notion of Gödel-numbering and view TNT on an entirely different level. It is not that
it is just difficult and complicated to write out the explanation on the TNT-level; it is
impossible. Such an explanation simply does not exist. There is, on the high level, a kind
of explanatory power which simply is lacking, in principle, on the TNT-level. G's
nontheoremhood is, so to speak, an intrinsically high-level fact. It is my suspicion that
this is the case for all undecidable propositions; that is to say: every undecidable
proposition is actually a Gödel sentence, asserting its own nontheoremhood in some
system via some code.

         Consciousness as an Intrinsically High-Level Phenomenon
Looked at this way, Gödel’s proof suggests-though by no means does it prove!-that there
could be some high-level way of viewing the mind/brain, involving concepts which do
not appear on lower levels, and that this level might have explanatory power that does not
exist-not even in principle-on lower levels. It would mean that some facts could be
explained on the high level quite easily, but not on lower levels at all. No matter how
long and cumbersome a low-level statement were made, it would not explain the
phenomena in question. It is the analogue to the fact that, if you make derivation after
derivation in TNT, no matter how long and cumbersome you make them, you will never
come up with one for G-despite the fact that on a higher level, you can see that G is true.
        What might such high-level concepts be? It has been proposed for eons, by
various holistically or "soulistically" inclined scientists and humanists, that consciousness
is a phenomenon that escapes explanation in terms of brain-components; so here is a
candidate, at least. There is also the ever-puzzling notion of free will. So perhaps these
qualities could be "emergent" in the sense of requiring explanations which cannot be
furnished by the physiology alone. But it is important to realize that if we are being
guided by Gödel’s proof in making such bold hypotheses, we must carry the




Strange Loops, Or Tangled Hierarchies                                                    703

analogy through thoroughly. In particular, it is vital to recall tnat is s nontheoremhood
does have an explanation-it is not a total mystery! The explanation- hinges on
understanding not just one level at a time, but the way in which one level mirrors its
metalevel, and the consequences of this mirroring. If our analogy is to hold, then,
"emergent" phenomena would become explicable in terms of a relationship between.
different levels in mental systems.,

                  Strange Loops as the Crux of Consciousness
My belief is that the explanations of "emergent" phenomena in our brains-for instance,
ideas, hopes, images, analogies, and finally consciousness and free will-are based on a
kind of Strange Loop, an interaction between levels in which the top level reaches back
down towards the bottom level and influences it, while at the same time being itself
determined by the bottom level. In other words, a self-reinforcing "resonance" between
different levels-quite like the Henkin sentence which, by merely asserting its own
provability, actually becomes provable. The self comes into being at the moment it has
the power to reflect itself.
        This should not be taken as an antireductionist position. It just implies that a
reductionistic explanation of a mind, in order to be comprehensible, must bring in "soft"
concepts such as levels, mappings, and meanings. In principle, I have no doubt that a
totally reductionistic but incomprehensible explanation of the brain exists; the problem is
how to translate it into a language we ourselves can fathom. Surely we don't want a
description in terms of positions and momenta of particles; we want a description which
relates neural activity to "signals" (intermediate-level phenomena)-and which relates
signals, in turn, to "symbols" and "subsystems", including the presumed-to-exist "self-
symbol". This act of translation from low-level physical hardware to high-level
psychological software is analogous to the translation of number-theoretical statements
into metamathematical statements. Recall that the level-crossing which takes place at this
exact translation point is what creates Godel's incompleteness and the self-proving
character of Henkin's sentence. I postulate that a similar level-crossing is what creates our
nearly unanalyzable feelings of self.
        In order to deal with the full richness of the brain/mind system, we will have to be
able to slip between levels comfortably. Moreover, we will have to admit various types of
"causality": ways in which an event at one level of description can "cause" events at other
levels to happen. Sometimes event A will be said to "cause" event B simply for the
reason that the one is a translation, on another level of description, of the other.
Sometimes "cause" will have its usual meaning: physical causality. Both types of
causality-and perhaps some more-will have to be admitted in any explanation of mind,
for we will have to admit causes that propagate both upwards and downwards in the
Tangled Hierarchy of mentality, just as in the Central Dogmap.




Strange Loops, Or Tangled Hierarchies                                                    704

        At the crux, then, of our understanding ourselves will come an understanding of
the Tangled Hierarchy of levels inside our minds. My position is rather similar to the
viewpoint put forth by the neuroscientist Roger Sperry in his excellent article "Mind,
Brain, and Humanist Values", from which I quote a little here:

   In my own hypothetical brain model, conscious awareness does get representation
   as a very real causal agent and rates an important place in the causal sequence and
   chain of control in brain events, in which it appears as an active, operational
   force.... To put it very simply, it comes down to the issue of who pushes whom
   around in the population of causal forces that occupy the cranium. It is a matter, in
   other words, of straightening out the peck-order hierarchy among intracranial
   control agents. There exists within the cranium a whole world of diverse causal
   forces; what is more, there are forces within forces within forces, as in no other
   cubic half-foot of universe that we know. ... To make a long story short, if one
   keeps climbing upward in the chain of command within the brain, one finds at the
   very top those over-all organizational forces and dynamic properties of the large
   patterns of cerebral excitation that are correlated with mental states or psychic
   activity.... Near the apex of this command system in the brain ... we find ideas. Man
   over the chimpanzee has ideas and ideals. In the brain model proposed here, the
   causal potency of an idea, or an ideal, becomes just as real as that of a molecule, a
   cell, or a nerve impulse. Ideas cause ideas and help evolve new ideas. They interact
   with each other and with other mental forces in the same brain, in neighboring
   brains, and, thanks to global communication, in far distant, foreign brains. And they
   also interact with the external surroundings to produce in toto a burstwise advance
   in evolution that is far beyond anything to hit the evolutionary scene yet, including
   the emergence of the living cell.'

        There is a famous breach between two languages of discourse: the subjective
language and the objective language. For instance, the "subjective" sensation of redness,
and the "objective" wavelength of red light. To many people, these seem to be forever
irreconcilable. I don't think so. No more than the two views of Escher's Drawing Hands
are irreconcilable from "in the system", where the hands draw each other, and from
outside, where Escher draws it all. The subjective feeling of redness comes from the
vortex of self-perception in the brain; the objective wavelength is how you see things
when you step back, outside of the system. Though no one of us will ever be able to step
back far enough to see the "big picture", we shouldn't forget that it exists. We should
remember that physical law is what makes it all happen-way, way down in neural nooks
and crannies which are too remote for us to reach with our high-level introspective
probes.

                          The Self-Symbol and Free Will
In Chapter XI I, it was suggested that what we call free will is a result of the interaction
between the self-symbol (or subsystem), and the other symbols in the brain. If we take
the idea that symbols are the high-level entities to


Strange Loops, Or Tangled Hierarchies                                                   705

which meanings should be attached, then we can' make a stab at explaining the
relationship between symbols, the self-symbol, and free will.
         One way to gain some perspective on the free-will question is to replace it by
what I believe is an equivalent question, but one which involves less loaded terms.
Instead of asking, "Does system X have free will?" we ask, "Does system X make
choices?" By carefully groping for what we really mean when we choose to describe a
system-mechanical or biological-as being capable of making "choices", I think we can
shed much light on free will it will be helpful to go over a few different systems which,
under various circumstances, we might feel tempted to describe as "making choices".
From these examples we can gain some perspective on what we really mean by the
phrase.
         Let us take the following systems as paradigms: a marble rolling down a bumpy
hill; a pocket calculator finding successive digits in the decimal expansion of the square
root of 2; a sophisticated program which plays a mean game of chess; a robot in a T-maze
(a maze with but a single fork, on one side of which there is a reward); and a human
being confronting a complex dilemma.
         First, what about that marble rolling down a hill? Does it make choices? I think
we would unanimously say that it doesn't, even though none of us could predict its path
for even a very short distance. We feel that it couldn't have gone any other way than it
did, and that it was just being shoved along by the relentless laws of nature. In our
chunked mental physics, of course, we can visualize many different "possible" pathways
for the marble, and we see it following only one of them in the real world. On some level
of our minds, therefore, we can't help feeling the marble has "chosen" a single pathway
out of those myriad mental ones; but on some other level of our minds, we have an
instinctive understanding that the mental physics is only an aid in our internal modeling
of the world, and that the mechanisms which make the real physical sequences of events
happen do not require nature to go through an analogous process of first manufacturing
variants in some hypothetical universe (the "brain of God") and then choosing between
them. So we shall not bestow the designation "choice" upon this process-although we
recognize that it is often pragmatically useful to use the word in cases like this, because
of its evocative power.
         Now what about the calculator programmed to find the digits of the square root of
2? What about the chess program? Here, we might say that we are just dealing with
"fancy marbles", rolling down "fancy hills". In fact, the arguments for no choice-making
here are, if anything, stronger than in the case of a marble. For if you attempt to repeat
the marble experiment, you will undoubtedly witness a totally different pathway being
traced down the hill, whereas if you rerun the square-root-of-2 program, you will get the
same results time after time. The marble seems to "choose" a different path each time, no
matter how accurately you try to reproduce the conditions of its original descent, whereas
the program runs down precisely the same channels each time.
         Now in the case of fancy chess programs, there are various possibilities.




Strange Loops, Or Tangled Hierarchies                                                  706

If you play a game against certain programs, and then start a second game with the same
moves as you made the first time, these programs will just move exactly as they did
before, without any appearance of having learned anything or having any desire for
variety. There are other programs which have randomizing devices that will give some
variety but not out of any deep desire. Such programs could be reset with the internal
random number generator as it was the first time, and once again, the same game would
ensue. Then there are other programs which do learn from their mistakes, and change
their strategy depending on the outcome of a game. Such programs would not play the
same game twice in a row. Of course, you could also turn the clock back by wiping out
all the changes in the memory which represent learning, just as you could reset the
random number generator, but that hardly seems like a friendly thing to do. Besides, is
there any reason to suspect that you would be able to change any of your own past
decisions if every last detail-and that includes your brain, of course-were reset to the way
it was the first time around?
         But let us return to the question of whether "choice" is an applicable term here. If
programs are just "fancy marbles rolling down fancy hills", do they make choices, or not?
Of course the answer must be a subjective one, but I would say that pretty much the same
considerations apply here as to the marble. However, I would have to add that the appeal
of using the word "choice", even if it is only a convenient and evocative shorthand,
becomes quite strong. The fact that a chess program looks ahead down the various
possible bifurcating paths, quite unlike a rolling marble, makes it seem much more like
an animate being than a square-root-of-2 program. However, there is still no deep self-
awareness here-and no sense of free will.
         Now let us go on to imagine a robot which has a repertoire of symbols. This robot
is placed in a T-maze. However, instead of going for the reward, it is preprogrammed to
go left whenever the next digit of the square root: of 2 is even, and to go right whenever it
is odd. Now this robot is capable of modeling the situation in its symbols, so it can watch
itself making choices. Each time the T is approached, if you were to address to the robot
the question, "Do you know which way you're going to turn this time?" it would have to
answer, "No". Then in order to progress, it would activate its "decider" subroutine, which
calculates the next digit of the square root of 2, and the decision is taken. However, the
internal mechanism of the decider is unknown to the robot-it is represented in the robot's
symbols merely as a black box which puts out "left"'s and "right"'s by some mysterious
and seemingly random rule. Unless the robot's symbols are capable of picking up the
hidden heartbeat of the square root of 2, beating in the L's and R's, it will stay baffled by
the "choices" which it is making. Now does this robot make choices? Put yourself in that
position. If you were trapped inside a marble rolling down a hill and were powerless to
affect its path, yet could observe it with all your human intellect, would you feel that the
marble's path involved choices? Of course not. Unless your mind is affecting the
outcome, it makes no difference that the symbols are present.




Strange Loops, Or Tangled Hierarchies                                                    707

So now we make a modification in our robot: we allow its symbols--including its self-
symbol-to affect the decision that is taken. Now here is an example of a program running
fully under physical law, which seems to get much more deeply at the essence of choice
than the previous examples did. When the robot's own chunked concept of itself enters
the scene, we begin to identify with the robot, for it sounds like the kind of thing we do. It
is no longer like the calculation of the square root of 2, where no symbols seem to be
monitoring the decisions taken. To be sure, if we were to look at the robot's program on a
very local level, it would look quite like the square-root program. Step after step is
executed, and in the end "left" or "right" is the output. But on a high level we can see the
fact that symbols are being used to model the situation and to affect the decision. That
radically affects our way of thinking about the program. At this stage, meaning has
entered this picture-the same kind of meaning as we manipulate with our own minds.

                     A Gödel Vortex Where All Levels Cross
Now if some outside agent suggests `L' as the next choice to the robot, the suggestion
will be picked up and channeled into the swirling mass of interacting symbols. There, it
will be sucked inexorably into interaction with the self-symbol, like a rowboat being
pulled into a whirlpool. That is the vortex of the system, where all levels cross. Here, the
`L' encounters a Tangled Hierarchy of symbols and is passed up and down the levels. The
self-symbol is incapable of monitoring all its internal processes, and so when the actual
decision emerges-'L' or 'R' or something outside the system-the system will not be able to
say where it came from. Unlike a standard chess program, which does not monitor itself
and consequently has no ideas about where its moves come from, this program does
monitor itself and does have ideas about its ideas-but it cannot monitor its own processes
in complete detail, and therefore has a sort of intuitive sense of its workings, without full
understanding. From this balance between self-knowledge and self-ignorance comes the
feeling of free will.
        Think, for instance, of a writer who is trying to convey certain ideas which to him
are contained in mental images. He isn't quite sure how those images fit together in his
mind, and he experiments around, expressing things first one way and then another, and
finally settles on some version. But does he know where it all came from? Only in a
vague sense. Much of the source, like an iceberg, is deep underwater, unseen-and he
knows that. Or think of a music composition program, something we discussed earlier,
asking when we would feel comfortable in calling it the composer rather than the tool of
a human composer. Probably we would feel comfortable when self-knowledge in terms
of symbols exists inside the program, and when the program has this delicate balance
between self-knowledge and self-ignorance. It is irrelevant whether the system is running
deterministically; what makes us call it a "choice maker" is whether we can identify with
a high-level description of the process which takes place when the




Strange Loops, Or Tangled Hierarchies                                                     708

             FIGURE 142. Print Gallery, by M. C. Escher (lithograph, 1956).

program runs. On a low (machine language) level, the program looks like any other
program; on a high (chunked) level, qualities such as "will", "intuition", "creativity", and
"consciousness" can emerge.
          The important idea is that this "vortex" of self is responsible for the tangledness,
for the Gödelian-ness, of the mental processes. People have said to me on occasion, "This
stuff with self-reference and so on is very amusing and enjoyable, but do you really think
  there is anything serious to it?" I certainly do. I think it will eventually turn out to be at
 the core of AI, and the focus of all attempts to understand how human minds work. And
              that is why Godel is so deeply woven into the fabric of my book.




Strange Loops, Or Tangled Hierarchies                                                      709

                   An Escher Vortex Where All Levels Cross
A strikingly beautiful, and yet at the same time disturbingly grotesque, illustration of the
cyclonic "eye" of a Tangled Hierarchy is given to us by Escher in his Print Gallery (Fig.
142). What we see is a picture gallery where a young man is standing, looking at a
picture of a ship in the harbor of a small town, perhaps a Maltese town, to guess from the
architecture, with its little turrets, occasional cupolas, and flat stone roofs, upon one of
which sits a boy, relaxing in the heat, while two floors below him a woman-perhaps his
mother-gazes out of the window from her apartment which sits directly above a picture
gallery where a young man is standing, looking at a picture of a ship in the harbor of a
small town, perhaps a Maltese town--What!? We are back on the same level as we began,
though all logic dictates that we cannot be. Let us draw a diagram of what we see (Fig.
143).




             FIGURE 143. Abstract diagram of M. C. Escher's Print Gallery.

What this diagram shows is three kinds of "in-ness". The gallery is physically in the town
("inclusion"); the town is artistically in the picture ("depiction"); the picture is mentally
in the person ("representation"). Now while this diagram may seem satisfying, in fact it is
arbitrary, for the number of levels shown is quite arbitrary. Look below at another way of
representing the top half alone (Fig. 144).




                FIGURE 144. A collapsed version of the previous figure.

Strange Loops, Or Tangled Hierarchies                                                    710

We have eliminated the "town" level; conceptually it was useful, but can just as well be
done without. Figure 144 looks just like the diagram for Drawing Hands: a Strange Loop
of two steps. The division markers are arbitrary, even if they seem natural to our minds.
This can be further accentuated by showing even more "collapsed" schematic diagrams of
Print Gallery, such as that in Figure 145.




                      FIGURE 145. Further collapse of Figure 143.

This exhibits the paradox of the picture in the starkest terms. Now-if the picture is "inside
itself", then is the young man also inside himself-, This question is answered in Figure
146.




                  FIGURE 146. Another way of collapsing Figure 143.

Thus, we see the young man "inside himself", in a funny sense which is made up of
compounding three distinct senses of "in".
        This diagram reminds us of the Epimenides paradox with its one-step self-
reference, while the two-step diagram resembles the sentence pair each of which refers to
the other. We cannot make the loop any tighter, but we can open it wider, by choosing to
insert any number of intermediate levels, such as "picture frame", "arcade", and
"building". If we do so, we will have many-step Strange Loops, whose diagrams are
isomorphic to those of Waterfall (Fig. 5) or Ascending and Descending (Fig. 6). The
number of levels is determined by what we feel is "natural", which may vary according to
context, purpose, or frame of mind. The Central Xmaps-Dog, Crab, Sloth, and Pipe-can
all be seen as involving three-step Strange Loops; alternatively, they can all be collapsed
into two- or one-step loops;. then again, they can be expanded out into multistage loops.
Where one perceives the levels is a matter of intuition and esthetic preference.
        Now are we, the observers of Print Gallery, also sucked into ourselves by virtue
of looking at it? Not really. We manage to escape that particular vortex by being outside
of the system. And when we look at the picture, we see things which the young man can
certainly not see, such as Escher's


Strange Loops, Or Tangled Hierarchies                                                    711

Signature, "MCE", in the central "blemish". Though the blemish seems like a defect,
perhaps the defect lies in our expectations, for in fact Escher could not have completed
that portion of the picture without being inconsistent with the rules by which he was
drawing the picture. That center of the whorl is-and must be-incomplete. Escher could
have made it arbitrarily small, but he could not have gotten rid of it. Thus we, on the
outside, can know that Print Gallery is essentially incomplete-a fact which the young
man, on the inside, can never know. Escher has thus given a pictorial parable for Gödel’s
Incompleteness Theorem. And that is why the strands of Gödel and Escher are so deeply
interwoven in my book.

                     A Bach Vortex Where All Levels Cross
One cannot help being reminded, when one looks at the diagrams of Strange Loops, of
the Endlessly Rising Canon from the Musical Offering. A diagram of it would consist of
six steps, as is shown in Figure 147. It is too .




   FIGURE 147. The hexagonal modulation scheme of Bach's Endlessly Rising Canon
               forms a true closed loop when Shepard tones are used.

bad that when it returns to C, it is an octave higher rather than at the exact original pitch.
Astonishingly enough, it is possible to arrange for it to return exactly to the starting pitch,
by using what are called Shepard tones, after the psychologist Roger Shepard, who
discovered the idea. The principle of a Shepard-tone scale is shown in Figure 148. In
words, it is this: you play parallel scales in several different octave ranges. Each note is
weighted independently, and as the notes rise, the weights shift. You make the top




Strange Loops, Or Tangled Hierarchies                                                      712

FIGURE 148: Two complete cycles of a Shephard tone scale, notated for piano. The
loudness of each note is proportional to its area, just as the top voice fades out, a new
bottom voice feebly enters. (Printed by Donald Boyd´s program “SMUT”.)


Strange Loops, Or Tangled Hierarchies                                                       713

octave gradually fade out, while at the same time you are gradually bringing in the
bottom octave. Just at the moment you would ordinarily be one octave higher, the
weights have shifted precisely so as to reproduce the starting pitch ... Thus you can go
"up and up forever", never getting any higher! You can try it at your piano. It works even
better if the pitches can be synthesized accurately under computer control. Then the
illusion is bewilderingly strong.
         This wonderful musical discovery allows the Endlessly Rising Canon to be played
in such a way that it joins back onto itself after going "up" an octave. This idea, which
Scott Kim and I conceived jointly, has been realized on tape, using a computer music
system. The effect is very subtle-but very real. It is quite interesting that Bach himself
was apparently aware, in some sense, of such scales, for in his music one can
occasionally find passages which roughly exploit the general principle of Shepard tones-
for instance, about halfway through the Fantasia from the Fantasia and Fugue in G Minor,
for organ.
         In his book J. S. Bach's Musical Offering, Hans Theodore David writes:

  Throughout the Musical Offering, the reader, performer, or listener is to search for
  the Royal theme in all its forms. The entire work, therefore, is a ricercar in the
  original, literal sense of the word.'

I think this is true; one cannot look deeply enough into the Musical Offering. There is
always more after one thinks one knows everything. For instance. towards the very end of
the Six-Part Ricercar, the one he declined to improvise, Bach slyly hid his own name,
split between two of the upper voices. Things are going on on many levels in the Musical
Offering. There are tricks with notes and letters; there are ingenious variations on the
King's Theme; there are original kinds of canons; there are extraordinarily complex
fugues; there is beauty and extreme depth of emotion; even an exultation in the many-
leveledness of the work comes through. The Musical Offering is a fugue of fugues, a
Tangled Hierarchy like those of Escher and Gödel, an intellectual construction which
reminds me, in ways I cannot express, of the beautiful many-voiced fugue of the human
mind. And that is why in my book the three strands of Gödel, Escher, and Bach are
woven into an Eternal Golden Braid.




Strange Loops, Or Tangled Hierarchies                                                 714

